A central idempotent lies in the centre $\Zc$, computed in Lemma \claimref{transport.centre}. If $n$ is even, $\Zc = \R$, whose idempotents are $0$ and $1$: no type C. If $n$ is odd, $\Zc = \R[\omega]$ with $\omega^2 = \pm 1$ by Lemma \claimref{transport.sign}. Case $\omega^2 = -1$: $\Zc \cong \C$, and an idempotent $a + b\omega$ satisfies $a^2 - b^2 = a$ and $2ab = b$; if $b \neq 0$ the second gives $a = \tfrac12$ and the first $-b^2 = \tfrac14$, impossible over $\R$; so $b = 0$ and $a \in \{0,1\}$: no type C. Case $\omega^2 = +1$: $\Zc \cong \R \oplus \R$ via $\varepsilon_\pm$, whose idempotents are $0, 1, \varepsilon_+, \varepsilon_-$; direct check $\varepsilon_\pm^2 = \tfrac14(1 \pm 2\omega + \omega^2) = \varepsilon_\pm$ and $\varepsilon_+\varepsilon_- = \tfrac14(1 - \omega^2) = 0$. By Lemma \claimref{transport.sign}, $n$ odd and $\omega^2 = +1$ happen together precisely when $p - q$ is odd and $p - q \equiv 0, 1 \pmod 4$, i.e.\ precisely when $p - q \equiv 1 \pmod 4$. For the exchange: for $n$ odd the pseudoscalar is central, so no inner automorphism, in particular no vector conjugation, moves it; the exchanging symmetry must instead be the \emph{grade involution}: $\hat{\omega} = (-1)^n \omega = -\omega$ for $n$ odd (\claimref{std-algebra.involution-properties}), so the grade involution, implemented on the algebra by conjugation in the larger algebra $\Cl(p,q{+}1)$ or realised abstractly as an outer automorphism, exchanges $\varepsilon_+ \leftrightarrow \varepsilon_-$. Within $A$ itself the two ideals are exchanged by no inner automorphism, centrality being conjugation-invariant, and this is exactly why type C is the deepest obstruction: no inner automorphism can act on it. The splitting into two simple ideals is Theorem \claimref{transport.classification-real-clifford-algebras} ($d \in \{1,5\}$), each ideal $\varepsilon_\pm A \cong \Mat_k(\mathbb{F})$.
